\chapter{AGI capability hypotheses and limits}
\label{evod51:chapter}
A system has memory, can call tools and can edit its own code. Which new
task can it now solve, under what budget, and how often does it fail? Without
those questions, a component inventory is easy to mistake for a capability
argument. Evolutor, DOGMA and Hermon DNA remain research directions; the
teaching implementations in this book do not establish artificial general
intelligence.

We turn a capability taxonomy and a claim ledger into operational contracts. Local tests, held-out demonstrations and independent replication
require different records. An executable metadata gate can reject an incomplete
claim, but it cannot determine whether an experimenter told the truth or whether
the chosen tests capture general intelligence.

\section{An evidence label is not a performance number}
\index{evidence level}\index{replication}\index{artifact identity}
\index{capability!evidence}
Use six evidence labels: speculative, hypothesized, implemented,
tested, benchmark-demonstrated and replicated. Speculative names a direction;
hypothesized adds a testable mechanism and falsifier. Implemented binds an
artifact. Tested adds local test evidence. Benchmark-demonstrated adds a bounded
empirical comparison. Replicated adds a separate reproduction of a specified
claim. These are editorial evidence categories, not a universally accepted
scientific maturity scale.

A serial promotion path is useful for preventing skipped checks, but passing
one edge cannot itself justify the next. A tested memory allocator remains
tested software even if its name contains ``intelligence.'' A hypothesized
capability can be falsified while the mechanism remains implemented. A new
scope or stronger wording requires a new claim rather than silently inheriting
old evidence.

\Plate{01-evidence}{Artifacts, local tests, bounded benchmark results and
independent reproduction support different edges. Each evidence node binds the
same claimed scope. A failure remains attached to that scope; adding a box to
the architecture does not supply an empirical edge.}{evod51:evidence}

The reference requires source and test digests for tested claims, then a
supporting evaluation with a baseline, data/artifact identity, declared budget,
sampling unit, held-out declaration, successful outcome and lower confidence
bound above a frozen threshold. Replication needs a distinct evidence ID,
a different evaluator and a declaration of independence under matching
suite/artifact/data/baseline/budget/unit/threshold. Here matching identities
define exact reproduction; robustness under changed data or hardware should
be a related, separately scoped claim.

\Listing{supports}
\Listing{promote}
These fields are necessary metadata only. A fake digest can have valid syntax,
an asserted held-out split can be contaminated, and two named teams can share
the same hidden development decisions. The code does not authenticate artifacts,
recompute statistics or inspect institutional independence. It deliberately
does not turn the sample ledger into evidence of an actual benchmark win.

\section{Breadth, depth and autonomy form a vector}
\index{autonomy}\index{generality}\index{capability vector}
The Levels of AGI framework separates performance and generality and examines
autonomy as an interaction/deployment choice \cite{evod51:levels}. We use
that distinction without assigning Evolutor a level. A broad assistant with
human intervention and a narrow unattended controller answer different
questions. Autonomy is not a multiplier that converts narrow skill into breadth.

Define a capability vector over declared task families, accompanied by budgets
and failure strata. A single mean can conceal a missing capability or a
catastrophic subgroup. The taxonomy below supplies operational tests, not a
claim that every term has a unique uncontested definition.

\subsection{Generalization, transfer and reasoning}
Generalization asks about held-out distributions, task rules or environments.
Name what is novel: examples, compositions, dynamics, goals or task families.
An unseen row generated by the same public template may test interpolation
rather than new-rule acquisition. Audit pretraining exposure, prompt examples,
retrieval corpora and human tuning against the final evaluation.

Transfer compares a source-trained or adapted system with a target-only control
under matched total opportunity. Report whether the source experience is charged
to the budget. A target score alone cannot show positive transfer. Reasoning
should be decomposed into compositional, symbolic, quantitative, spatial,
causal or interactive tasks with checkable outputs. Plausible intermediate
prose is not proof of a valid latent mechanism. Use answer verifiers, minimal
counterexamples and interventions on task structure rather than scoring only
the fluency of explanations.

\subsection{Planning, world modeling and tool use}
Planning requires goal achievement through actions under cost and uncertainty.
Measure success, actions, violations, recovery, wall time and human assistance.
A static plan-shaped answer has not yet interacted with a changing environment.
World modeling asks whether predicted transitions or latent dynamics improve
held-out decisions. Compare against reactive, retrieval-only and model-free
controls while fixing action and compute budgets.

Tool use requires correct selection, typed arguments, error recovery and task
completion under deterministic authorization. A connector proves availability,
not competence. Remove a needed tool, introduce a controlled timeout or return
a version conflict. The outcome should distinguish justified deferral from
invalid execution, not reward any completed action regardless of policy.

\subsection{Memory, continual learning and adaptation}
Memory evaluation names the storage type: transient carry, context, addressable
records or learned parameters. Measure retention horizon, retrieval precision
and recall, interference, provenance and deletion. A record surviving on disk
does not imply the model can retrieve or use it. Chapter 50's acquisition and
retention matrix is the corresponding contract for continual learning; saving
events alone is not weight learning.

Adaptation declares which of context, memory, weights, structure or population
changed. Freeze data access, number of trials and update budget. A task-ID
oracle changes the problem relative to a task-free stream. Performance after
hundreds of target-specific trials is not evidence of efficient acquisition
from a few observations.

\subsection{Metacognition, self-improvement and open-endedness}
Operational metacognition includes calibrated uncertainty, detection of errors,
strategy switching and deferral when evidence is insufficient. Evaluate selective
risk at several coverage levels: a system that refuses everything has low
observed error but no useful coverage. Confidence text without probability/error
calibration is weak evidence.

Self-editing produces a candidate change; self-improvement requires repeated,
causally attributable held-out gains with declared total cost and regression
limits. Compare against equally budgeted random edits, human-designed updates
or fixed search. Preserve negative generations. Deployment authority stays
outside the proposal mechanism. Open-ended learning additionally requires
sustained useful novelty beyond a fixed suite, with task/niche definitions,
novelty measures, transfer and resource accounting. New strings or new filenames
do not establish new skills.

\section{Unknown-task evaluation needs an information firewall}
\index{held-out evaluation}\index{benchmark contamination}\index{interactive evaluation}
Before a final run, freeze artifact, prompts, sampler, retrieval corpus,
adaptation algorithm, baseline and allowed tools. Separate development tasks,
adaptation observations and final evaluation. The evaluator may reveal the
observations a new environment legitimately supplies without exposing withheld
answers. Task selection and human interventions are part of the protocol.

\Plate{02-holdout}{A learner sees permitted interaction observations but no
final labels. The evaluator scores sealed trajectories after artifact and
adaptation rules are frozen. Repeated development feedback from the final suite
breaks this boundary even if the filenames still say test.}{evod51:holdout}

ARC-AGI-2 and ARC-AGI-3 provide specific benchmark designs, not universal
certificates. The official ARC-AGI-2 page emphasizes novel reasoning and
efficiency reporting \cite{evod51:arc2}. The ARC-AGI-3 page describes
interactive environments, exploration and action efficiency
\cite{evod51:arc3}. Their versions and access dates belong in a result
capsule because benchmark policies evolve. We import no benchmark scores and
make no claim that Evolutor ran either suite.

Interactive evaluation must distinguish environmental reset, episode context,
cross-episode memory and parameter learning. If task instances share a secret
rule, carrying memory across instances changes the independence assumptions.
Charge environment actions, tokens, model calls and any human help separately.
Report failed runs and budget exhaustion as outcomes rather than silently
dropping them from the denominator.

\section{Skill acquisition has a budget axis, not a magical ratio}
\index{skill acquisition}\index{budget axis}\index{acquisition area}
A conceptual ratio such as $\Delta\text{capability}/(\text{data}+\text{compute}
+\text{actions}+\text{risk})$ identifies an important concern but adds unlike
units. Do not print it as an empirical efficiency metric without a declared
conversion policy. Use separate curves or a vector of constraints instead.

For one chosen action axis $0=b_0<\cdots<b_K=B$, let $q_k\in[0,1]$ be
quality after budget $b_k$. Under linear interpolation, a normalized area is
\begin{equation}\label{evod51:area}
\overline q_B=\frac1B\sum_{k=1}^K
(b_k-b_{k-1})\frac{q_k+q_{k-1}}2.
\end{equation}
It lies in $[0,1]$ because it averages bounded interpolated quality over the
budget interval. It changes if one chooses a different interpolation or budget
axis. It cannot compare dollars with actions without an explicit conversion.
\Listing{acquisition_area}
\begin{center}\input{\Generated/areas.tex}\end{center}
Both illustrative curves reach one at budget two. The early curve $(0,.8,1)$
has area $.65$; the late curve $(0,.2,1)$ has area $.35$. Final score alone
loses this distinction. Multiplying all budget coordinates by a positive constant
preserves the normalized area, a tested unit-rescaling property.

\Plate{03-acquisition}{The same final score conceals different acquisition
histories. Areas assume linear interpolation between the three stated points.
The curves are assigned numbers, not results from an agent benchmark.}
{evod51:acquisition}

\section{Finite tests cannot prove unrestricted generality}
\index{finite observational underdetermination}\index{Hoeffding bound}\index{zero-failure bound}
\begin{proposition}[Finite observational underdetermination]
Suppose an evaluator observes a finite set of task histories and the system
class permits arbitrary behavior on an unobserved history. Two systems can
agree on every observed output while differing on that unobserved history.
\end{proposition}
\begin{proof}
Take one system $f$. Construct $g$ to return $f$ on every observed history
and a different output on an unobserved history $h^*$. By assumption the class
permits that construction. The finite transcript is identical for $f$ and $g$,
so it cannot establish their agreement at $h^*$.
\end{proof}
The conclusion is limited: evidence supports a distributional or structural
claim under assumptions, not universal competence. It does not make benchmarks
useless. Hidden tasks, broad coverage, strong priors and mechanistic arguments
can strengthen the inference while leaving its scope explicit.

For independent bounded task-family success values in $[0,1]$, Hoeffding's
inequality gives $\Pr(|\hat q-q|\ge\epsilon)\le2e^{-2n\epsilon^2}$.
For $m$ frozen claims, a union bound gives total error at most
$2m e^{-2n\epsilon^2}$. Thus $n\ge\log(2m/\delta)/(2\epsilon^2)$ is a
sufficient sample count under those assumptions. Correlated puzzles, adaptive
selection and repeated inspection invalidate the simple independent-sample
interpretation. The units must be task families if that is the claimed population,
not thousands of near-duplicate instances treated as independent breadth.

Zero observed violations is also weaker than zero risk. For $n$ independent
Bernoulli trials with constant failure probability $p$, zero failures has
probability $(1-p)^n$. Solving $(1-p)^n=\alpha$ gives a one-sided upper bound
\[
p_{\mathrm{upper}}=1-\alpha^{1/n}.
\]
\Listing{zero_failure_upper}
\begin{center}\input{\Generated/zero_failures.tex}\end{center}
These generated bounds are a mathematical illustration, not a certification
of an agent's safety. Dependence, adversarial distribution shift and missing
failure detection can make the premise inappropriate.

\section{Three research directions, three falsifiable hypotheses}
\index{DOGMA!hypothesis}\index{Hermon DNA!hypothesis}\index{self-improvement}
DOGMA's non-Transformer state organization motivates quality--memory--latency
comparisons with strong recurrent/SSM and other appropriate sequence baselines.
State boundedness alone says nothing about retained task information. If matched
long-context tasks show no advantage within the predeclared effect margin,
reject that bounded efficiency claim rather than changing the task list afterward.
Chapter 31's diagonal recurrence is narrower than full DOGMA and cannot certify
that comparison.

Hermon DNA preserves Transformer semantics while testing genomic structure and
serving machinery. A task-trained genomic improvement can establish domain skill;
it does not establish broad reasoning. Compare a DNA-aware mechanism with an
otherwise matched Transformer, and separate serving improvements from learned
quality. Chapter 35's learned-absolute-position cache is a semantic implementation,
not a new pretrained foundation model.

Evolutor's compiler, memory and governance direction motivates persistent
workflow tests: can a frozen system recover, use provenance-aware memory and
complete held-out tasks within policy and budget? Cross-engine integration is
an engineering property until those evaluations run. The book's runtime,
memory and authority references are separately bounded contracts; their
composition into a durable service and its task evaluations are not assumed
by this chapter's ledger.

Capability and deployment safety remain separate axes. A capable model can
violate policy; a correctly fenced tool host can serve a limited model. Measure
policy violations, isolation, revocation and recovery alongside success. More
human intervention may improve safety and success while reducing autonomy;
report all three rather than hiding assistance in the runtime.

\section{Selective glossary}
\begin{description}
\item[Capability scope.] The task population, information access and budget
to which a behavioral claim applies.
\item[Held-out evaluation.] Evidence withheld from development and selection,
with permitted online adaptation observations separately specified.
\item[Replication.] A distinct evaluation of the same bounded claim by a
separate evaluator; metadata declarations alone do not verify independence.
\item[Acquisition area.] A budget-axis-dependent integral of quality under a
declared interpolation; final score alone does not determine it.
\item[Autonomy.] The degree of independent task execution and recovery relative
to human assistance, distinct from breadth or performance.
\end{description}

\section{Exercises and worked solutions}
\begin{enumerate}
\item A planner object exists. What evidence level follows? \textbf{Answer:}
implemented if artifact-bound; successful planning still needs task evidence.
\item Memory survives restart. Does continual learning follow? \textbf{Answer:}
no; acquisition, retention and the update layer remain untested.
\item Why require a falsifier? \textbf{Answer:} it names a result that would
weaken or reject the claim instead of allowing every outcome to count as success.
\item Can one evidence record establish replication? \textbf{Answer:} no;
the ledger requires a distinct evaluator/record under a matching contract.
\item Is declared independence proof of independence? \textbf{Answer:} no;
it is metadata requiring external verification.
\item Why cannot promotion broaden scope? \textbf{Answer:} evidence for one
task population does not automatically support a larger population.
\item A tool connector is available. What remains? \textbf{Answer:} selection,
arguments, recovery, authorized completion and held-out success under budget.
\item Compare autonomy with breadth. \textbf{Answer:} intervention independence
and task coverage are different dimensions.
\item Does fluent reasoning text prove causal reasoning? \textbf{Answer:}
no; use checkable outcomes and task-structure interventions.
\item Why is final-suite feedback dangerous? \textbf{Answer:} repeated tuning
turns the supposedly held-out suite into development information.
\item Compute the early curve area. \textbf{Answer:} $(.4+.9)/2=.65$.
\item Compute the late curve area. \textbf{Answer:} $(.1+.6)/2=.35$.
\item Rescale actions by ten. \textbf{Answer:} both numerator and denominator
scale by ten, preserving normalized area.
\item Why not sum actions and risk in a denominator? \textbf{Answer:}
different units need explicit policy conversions; a vector is often clearer.
\item Construct the finite-test counterexample. \textbf{Answer:} copy all
tested outputs and change one unobserved history; the transcript stays identical.
\item Derive zero-failure upper bound. \textbf{Answer:} invert
$(1-p)^n=\alpha$, giving $1-\alpha^{1/n}$.
\item Does zero failure among correlated tasks use the same bound? \textbf{Answer:}
not under the stated independent Bernoulli premise.
\item Design self-improvement evaluation. \textbf{Rubric:} external frozen
evaluator, matched search control, held-out gains across iterations, charged
total budget, regressions, lineage and separately authorized promotion.
\item Design an interactive breadth test. \textbf{Rubric:} distinct held-out
task families, observation/action rules, assistance accounting, contamination
audit, separate success, cost, calibration and failure metrics, and task-level uncertainty.
\item Strongest conclusion here? \textbf{Answer:} executable evidence-contract
checks and falsifiable research hypotheses; no demonstrated AGI capability.
\end{enumerate}

\section{Bridge}
Chapter 52 binds claims to experiments, failures and artifact closures. A final
synthesis must make unresolved work easier to find, including evidence that
contradicts the preferred architecture.
